\documentclass[11pt]{article}
\usepackage[T1]{fontenc}
\usepackage{textcomp}
\usepackage[margin=0.75in]{geometry}
\usepackage{amsmath,amssymb,amsthm}
\usepackage{array,booktabs}
\usepackage{hyperref}
\setlength{\parindent}{0pt}
\newtheorem{theorem}{Theorem}
\theoremstyle{definition}
\newtheorem{definition}{Definition}
\title{Flow Proof Artifact: prop-01-triangles-parallelograms-same-height}
\author{Generated by \texttt{flow doc proof}}
\date{Flow Proof Book}
\begin{document}
\maketitle
Triangles and parallelograms under the same height are to one another as their bases.\\[0.5em]
\textbf{Source.} euclid — Elements, Book VI, Proposition 1\\[0.5em]
\bigskip
\noindent{\large \textbf{Derived fact 1.} \textit{Triangles and parallelograms under the same height are to one another as their bases} \hfill the Euclidean plane \cdot Euclid Book VI \cdot Proposition 1: triangles and parallelograms of the same height are as their bases}\par
\smallskip\noindent\textit{Coordinate: the Euclidean plane \cdot Euclid Book VI \cdot Proposition 1: triangles and parallelograms of the same height are as their bases}\par
\smallskip\noindent\textit{Source: euclid — Elements, Book VI, Proposition 1}\par
\smallskip\noindent\textit{Built on: proposition 1: equimultiples of equimultiples are equimultiple of the originals, for Euclid Book V on the Euclidean plane}\par
\medskip
\noindent\textbf{Goal.} Triangles and parallelograms under the same height are to one another as their bases\par
\noindent$\displaystyle \text{area ratio equals base ratio}$\par
\medskip
\noindent
\begin{tabular}{@{} >{\raggedright\arraybackslash}p{0.06\textwidth} >{\raggedright\arraybackslash}p{0.47\textwidth} >{\raggedleft\arraybackslash}p{0.41\textwidth} @{}}
\textbf{\#} & \textbf{Proof} & \textbf{Mathematics} \\
\midrule
\textbf{1.} & We prove that proposition 1: triangles and parallelograms of the same height are as their bases for Euclid Book VI the Euclidean plane. & \\[0.45em]
\textbf{2.} & Let ABC = triangle. & \\[0.45em]
\textbf{3.} & Let base\_BC = base. & \\[0.45em]
\textbf{4.} & Let height = common\_height. & \\[0.45em]
\textbf{5.} & From step 2, step 3, and step 4, we can deduce that area(ABC) : area(DEF) equals base BC : base EF. & $\displaystyle area(ABC) : area(DEF) = \text{base BC} : \text{base EF}$ \\[0.45em]
\textbf{6.} & We invoke the derived fact governing Euclid Book V the Euclidean plane: proposition 1: equimultiples of equimultiples are equimultiple of the originals, for Euclid Book V on the Euclidean plane. & \\[0.45em]
\textbf{7.} & From step 6, this implies area ratio equals base ratio. Hence proven. & $\displaystyle \text{area ratio equals base ratio}$ \\[0.45em]
\bottomrule
\end{tabular}
\label{thm:Geometry-Euclid Book VI-proposition_1_triangles_and_parallelograms_of_the_same_height_are_as_their_bases}
\par\medskip\hrule
\end{document}
